% GENERATED FILE. Edit canonical lessons, then run npm run generate:books.
% canonical-source-hash: fnv1a64:06432e806b11f775
% canonical-lessons: GE-C179-empfangen, GE-C179-erinnern, GE-C179-erlauben, GE-C179-ernaehren, GE-C179-erreichen

\chapter{To receive, To remind, To allow, To feed, To reach}
\label{ch:ge-to-receive-to-remind-to-allow-to-feed-to-reach}
\hlchaptermodality{179}

\begin{chapteropening}
\textbf{By the end of this chapter:} \emph{I can say to receive, to remind, to allow, to feed and to reach.}

Complete the last lesson of chapter 179: I can say to receive, to remind, to allow, to feed and to reach.
\end{chapteropening}

\section[empfangen]{empfangen — to receive, to welcome}

\label{lesson:GE-C179-empfangen}



\begin{quote}
Before the new one: say the German for to get off, then the German for to change.
\end{quote}

\subsection*{You'll want to know: empfangen}
\textbf{empfangen} — \textquotedblleft{}to receive, to welcome\textquotedblright{}.

\begin{grammarlens}[title={What you've built}]
One more word for this chapter.
\end{grammarlens}

\subsection*{Your turn}
\begin{itemize}
\raggedright
  \item \emph{Say it:} \emph{empfangen}
  \item \emph{Say it:} \emph{empfangen}, once more
  \item \emph{Say it:} say \emph{empfangen}
  \item \emph{Recall:} say the German for to get off, then the German for to change, then say \emph{empfangen} again
\end{itemize}

\begin{tcolorbox}[breakable,colback=teal!4,colframe=teal!35!black,title={Before you move on}]
What does empfangen mean? (\textquotedblleft{}To receive, to welcome\textquotedblright{}.) Say it once more.
\end{tcolorbox}
\section[erinnern]{erinnern — to remind (sich erinnern: to remember)}

\label{lesson:GE-C179-erinnern}



\begin{quote}
Before the new one: say the German for to change, then the German for to receive.
\end{quote}

\subsection*{You'll want to know: erinnern}
\textbf{erinnern} — \textquotedblleft{}to remind (sich erinnern: to remember)\textquotedblright{}.

\begin{grammarlens}[title={What you've built}]
One more word for this chapter.
\end{grammarlens}

\subsection*{Your turn}
\begin{itemize}
\raggedright
  \item \emph{Say it:} \emph{erinnern}
  \item \emph{Say it:} \emph{erinnern}, once more
  \item \emph{Say it:} say \emph{erinnern}
  \item \emph{Recall:} say the German for to change, then the German for to receive, then say \emph{erinnern} again
\end{itemize}

\begin{tcolorbox}[breakable,colback=teal!4,colframe=teal!35!black,title={Before you move on}]
What does erinnern mean? (\textquotedblleft{}To remind (sich erinnern: to remember)\textquotedblright{}.) Say it once more.
\end{tcolorbox}
\section[erlauben]{erlauben — to allow}

\label{lesson:GE-C179-erlauben}



\begin{quote}
Before the new one: say the German for to receive, then the German for to remind.
\end{quote}

\subsection*{You'll want to know: erlauben}
\textbf{erlauben} — \textquotedblleft{}to allow\textquotedblright{}.

\begin{grammarlens}[title={What you've built}]
One more word for this chapter.
\end{grammarlens}

\subsection*{Your turn}
\begin{itemize}
\raggedright
  \item \emph{Say it:} \emph{erlauben}
  \item \emph{Say it:} \emph{erlauben}, once more
  \item \emph{Say it:} say \emph{erlauben}
  \item \emph{Recall:} say the German for to receive, then the German for to remind, then say \emph{erlauben} again
\end{itemize}

\begin{tcolorbox}[breakable,colback=teal!4,colframe=teal!35!black,title={Before you move on}]
What does erlauben mean? (\textquotedblleft{}To allow\textquotedblright{}.) Say it once more.
\end{tcolorbox}
\section[ernähren]{ernähren — to feed, to nourish}

\label{lesson:GE-C179-ernaehren}



\begin{quote}
Before the new one: say the German for to remind, then the German for to allow.
\end{quote}

\subsection*{You'll want to know: ernähren}
\textbf{ernähren} — \textquotedblleft{}to feed, to nourish\textquotedblright{}.

\begin{grammarlens}[title={What you've built}]
One more word for this chapter.
\end{grammarlens}

\subsection*{Your turn}
\begin{itemize}
\raggedright
  \item \emph{Say it:} \emph{ernähren}
  \item \emph{Say it:} \emph{ernähren}, once more
  \item \emph{Say it:} say \emph{ernähren}
  \item \emph{Recall:} say the German for to remind, then the German for to allow, then say \emph{ernähren} again
\end{itemize}

\begin{tcolorbox}[breakable,colback=teal!4,colframe=teal!35!black,title={Before you move on}]
What does ernähren mean? (\textquotedblleft{}To feed, to nourish\textquotedblright{}.) Say it once more.
\end{tcolorbox}
\section[erreichen]{erreichen — to reach}

\label{lesson:GE-C179-erreichen}



\begin{quote}
Before the new one: say the German for to allow, then the German for to feed.
\end{quote}

\subsection*{You'll want to know: erreichen}
\textbf{erreichen} — \textquotedblleft{}to reach\textquotedblright{}.

\begin{grammarlens}[title={What you've built}]
One more word for this chapter.
\end{grammarlens}

\subsection*{Your turn}
\begin{itemize}
\raggedright
  \item \emph{Say it:} \emph{erreichen}
  \item \emph{Say it:} \emph{erreichen}, once more
  \item \emph{Say it:} say \emph{erreichen}
  \item \emph{Recall:} say the German for to allow, then the German for to feed, then say \emph{erreichen} again
\end{itemize}

\begin{tcolorbox}[breakable,colback=teal!4,colframe=teal!35!black,title={Before you move on}]
What does erreichen mean? (\textquotedblleft{}To reach\textquotedblright{}.) Say it once more.
\end{tcolorbox}
